\documentclass[11pt]{article}

\usepackage[margin=1in]{geometry}
\usepackage{amsmath,amssymb,amsthm}
\usepackage{booktabs}
\usepackage{enumitem}
\usepackage{hyperref}
\usepackage{parskip}

\hypersetup{
  colorlinks=true,
  linkcolor=blue,
  urlcolor=blue,
}

% Theorem environments
\newtheorem{theorem}{Theorem}
\newtheorem{lemma}[theorem]{Lemma}
\theoremstyle{definition}
\newtheorem{definition}{Definition}

% Shortcuts
\newcommand{\joins}{\bowtie}
\newcommand{\fkjoin}[1]{\joins\![#1]}
\newcommand{\set}[1]{\{#1\}}
\newcommand{\Uf}{U_f}
\newcommand{\Up}{U_p}
\newcommand{\Rf}{R_f}
\newcommand{\Rp}{R_p}
\newcommand{\Nf}{N_f}
\newcommand{\Np}{N_p}
\newcommand{\Cf}{C_f}
\newcommand{\Cp}{C_p}
\newcommand{\Jf}{J_f}
\newcommand{\Jp}{J_p}
\newcommand{\outerf}{\mathrm{outer}_f}
\newcommand{\outerp}{\mathrm{outer}_p}
\newcommand{\nn}{\mathrm{nn}}
\newcommand{\unique}{\mathrm{unique}}
\newcommand{\anchor}{\mathrm{anchor}}
\newcommand{\cols}{\mathrm{cols}}
\newcommand{\FK}{\mathrm{FK}}

\title{Foreign Key Joins --- Mathematical Specification}
\author{}
\date{}

\begin{document}
\maketitle

\tableofcontents
\bigskip

%% ====================================================================
\section{Notation}

Relations are multisets (bags) of tuples.  We write $\pi_X(M)$ for
the multiset projection of $M$ onto columns $X$: each tuple in $M$
is restricted to the columns in~$X$, preserving multiplicities.  The
join of relations $R$ and $S$ on predicate~$\theta$ is written
$R \joins_\theta S$.


%% ====================================================================
\section{Foreign Key Join}

\begin{definition}[Foreign Key Join]
A \emph{foreign key join} combines a referencing relation~$R$ with a
referenced relation~$S$ on column lists
$\mathbf{a} \subseteq \cols(R)$ and
$\mathbf{b} \subseteq \cols(S)$.
\end{definition}

\paragraph{Traceability.}
Every column in $\mathbf{a}$ traces---by projection and renaming
only, no expressions or aggregation---to a column of a single base
table~$r$.  Likewise every column in $\mathbf{b}$ traces to a single
base table~$s$.  Write $\mathbf{a'}$, $\mathbf{b'}$ for the
corresponding base-table columns.

\paragraph{Validity.}
The foreign key join is \emph{valid} if and only if:
\begin{enumerate}
\item A referential constraint
  $\FK(r,\;\mathbf{a'},\;s,\;\mathbf{b'})$ exists.
\item Multiset equality\footnote{Two multisets are equal if and only
  if every element appears with the same multiplicity, including
  elements containing \textsc{null}.} holds:
  \[
    \pi_{\mathbf{b}}(S) \;=\; \pi_{\mathbf{b'}}(s)
  \]
\end{enumerate}

\paragraph{Result.}
A valid foreign key join produces
\[
  R \joins_{\mathbf{a}=\mathbf{b}} S
\]
under any join type
$\tau \in \{\text{inner},\;\text{left},\;\text{right},\;\text{full}\}$.

Condition~2 is trivially satisfied when $S$ is the base table~$s$
itself.  For derived relations (views, CTEs, subqueries), the database
must verify it statically from the schema and query structure.


%% ====================================================================
\section{Static Verification}

For base tables, Condition~2 holds trivially:
$\pi_{\mathbf{b}}(S) = \pi_{\mathbf{b'}}(s)$ when $S = s$ and
$\mathbf{b} = \mathbf{b'}$.

For derived relations, the database cannot evaluate $\pi_{\mathbf{b}}(S)$
at parse time---it depends on the data.  Instead, it must verify
Condition~2 from the query structure alone.

We define a static analysis that is \emph{sound} but not
\emph{complete}: if it accepts, Condition~2 is guaranteed to hold
(soundness); but it may reject some queries for which Condition~2 does
in fact hold (incompleteness).  Soundness means no false positives
(accepting invalid joins); incompleteness means some false negatives
(rejecting valid ones).

The approach: track four properties---uniqueness, row preservation,
null preservation, and chain reachability---bottom-up through the join
tree, then check three conditions at the entry point.


%% ====================================================================
\section{Join Trees}

\subsection{Nodes}

A join tree is built from:

\begin{itemize}
\item \textbf{Leaf nodes}: a base table~$t$, optionally marked as
  \emph{filtered}.  A leaf is filtered when the base table is
  accessed through a derived relation whose query restricts rows
  (through selection, limit, or similar mechanisms).  An unfiltered
  leaf accesses the complete base table.

\item \textbf{FK join nodes}:
  $J = (\Jf \fkjoin{f \to p,\;\nn,\;\tau} \Jp)$ with parameters:
  \begin{itemize}[nosep]
  \item $\Jf$: the subtree containing the referencing relation
  \item $\Jp$: the subtree containing the referenced relation
  \item $f$: the base table identifier of the referencing table
  \item $p$: the base table identifier of the referenced table
  \item $\nn$: a boolean---true iff the FK columns carry
    \textsc{not null} constraints
  \item $\tau \in
    \{\text{inner},\;\text{left},\;\text{right},\;\text{full}\}$:
    the join type
  \end{itemize}
\end{itemize}

Only FK joins appear as join nodes.  Non-FK joins inside a derived
relation prevent analysis---the algorithm cannot reason about their
row-preservation properties, so their output sets are empty.


\subsection{Outer Preservation}

Define the predicate $\mathrm{outer}(\mathit{side}, \tau)$ as true
when the given side is on the outer (preserved) side of the join.  In
a FK join $J = (\Jf \joins \Jp)$, ``left'' corresponds to~$\Jf$
(referencing) and ``right'' corresponds to~$\Jp$ (referenced):

\begin{center}
\begin{tabular}{lcc}
\toprule
$\tau$ & $\outerf$ (referencing preserved)
       & $\outerp$ (referenced preserved) \\
\midrule
inner & false & false \\
left  & true  & false \\
right & false & true  \\
full  & true  & true  \\
\bottomrule
\end{tabular}
\end{center}

When a side is preserved, all its rows appear in the output.  Rows
without a match on the other side are padded with \textsc{null}s.


%% ====================================================================
\section{Tracked Properties}

Each node~$J$ carries four properties computed bottom-up over the
universe~$T$ of base table identifiers in $J$'s subtree:

\begin{center}
\begin{tabular}{clcl}
\toprule
Symbol & Name & Type & Meaning \\
\midrule
$U$ & uniqueness
    & $\subseteq T$
    & Tables whose unique-key uniqueness is preserved \\
$R$ & row preservation
    & $\subseteq T$
    & Tables whose complete row set appears in $J$'s result \\
$N$ & null preservation
    & $\subseteq T$
    & Tables whose key-column values are intact \\
$C$ & chains
    & $\subseteq T \times T$
    & $(A,B) \in C$ means $A$ reaches $B$ via a NOT NULL FK path \\
\bottomrule
\end{tabular}
\end{center}

\paragraph{Chain invariant.}
$(A, B) \in C$ if and only if there exists a directed path
$A = t_0 \to t_1 \to \cdots \to t_k = B$ where:
\begin{enumerate}
\item $A \in R$ \quad ($A$'s complete row set appears in $J$'s result)
\item Each edge $t_i \to t_{i+1}$ corresponds to a NOT NULL FK join
  with $t_i$ as the referencing table and $t_{i+1}$ as the referenced
  table
\item At the time edge $t_i \to t_{i+1}$ was processed, $t_{i+1}$
  was row-preserving ($t_{i+1} \in R$ of the referenced subtree at
  that node)
\end{enumerate}


%% ====================================================================
\section{Base Cases}

At a leaf node for base table~$t$:

\paragraph{Unfiltered.}
\[
  U = \set{t},\quad R = \set{t},\quad N = \set{t},\quad C = \emptyset
\]

\paragraph{Filtered.}
\[
  U = \set{t},\quad R = \emptyset,\quad N = \set{t},\quad C = \emptyset
\]

Uniqueness is always preserved: filtering cannot create duplicate
keys.  Row preservation is lost because filtering may exclude rows.
Null preservation holds because no join padding has occurred.


%% ====================================================================
\section{Propagation Rules}

Given $J = (\Jf \fkjoin{f \to p,\;\nn,\;\tau} \Jp)$ with input
properties $(\Uf, \Rf, \Nf, \Cf)$ from the referencing subtree and
$(\Up, \Rp, \Np, \Cp)$ from the referenced subtree, compute the
output properties $(U', R', N', C')$.


\subsection{Outer Preservation Predicate}

The shorthand $\outerf$ and $\outerp$ denotes whether the
referencing and referenced sides are outer-preserved, as defined in
\S4.2:

\begin{center}
\begin{tabular}{lcc}
\toprule
$\tau$ & $\outerf$ & $\outerp$ \\
\midrule
inner & false & false \\
left  & true  & false \\
right & false & true  \\
full  & true  & true  \\
\bottomrule
\end{tabular}
\end{center}


\subsection{Uniqueness (\texorpdfstring{$U$}{U})}

\[
  U' = \Uf \;\cup\;
    \begin{cases}
      \Up & \text{if } \unique(f) \land f \in \Uf \\
      \emptyset & \text{otherwise}
    \end{cases}
\]
where $\unique(f)$ means the FK columns form a unique key on~$f$.

The referencing side's uniqueness set is always inherited.  The
referenced side's set is inherited only when the FK columns are unique
and $f$ itself preserves uniqueness---making the join one-to-one, so
no referenced rows are duplicated.


\subsection{Null Preservation (\texorpdfstring{$N$}{N})}

Start with the union of both inputs:
\[
  N' = \Nf \cup \Np
\]

Then remove entries from the null-padded (inner) side:

\begin{itemize}
\item If $\outerf \land \lnot\nn$: set $N' = N' \setminus \Np$.

  When the referencing side is preserved, unmatched referencing rows
  produce ghost rows with \textsc{null}s on the referenced side.  All
  tables from $\Jp$ have their column values corrupted in these ghost
  rows.

  Exception: if $\nn$ holds, the FK guarantee ensures every
  referencing row matches a referenced row (the FK value is non-null
  and the constraint ensures a match exists).  No unmatched
  referencing rows exist, so no ghost rows appear and $\Np$ is not
  removed.

\item If $\outerp$: set $N' = N' \setminus \Nf$.

  When the referenced side is preserved, unmatched referenced rows
  produce ghost rows with \textsc{null}s on the referencing side.
  All tables from $\Jf$ have their column values corrupted.  This
  removal is unconditional---the FK constraint provides no guarantee
  in this direction (not every referenced row need be referenced).
\end{itemize}

\paragraph{Net effect on entries from each side:}

\begin{center}
\begin{tabular}{lll}
\toprule
$\tau$ & $\Nf$ entries & $\Np$ entries \\
\midrule
inner & inherited & inherited \\
left  & inherited & removed if $\lnot\nn$;\; inherited if $\nn$ \\
right & removed   & inherited \\
full  & removed   & removed if $\lnot\nn$;\; inherited if $\nn$ \\
\bottomrule
\end{tabular}
\end{center}

Here ``inherited'' means the entries pass through from the input, and
``removed'' means they are unconditionally absent from the output.


\subsection{Row Preservation and Chains (\texorpdfstring{$R$, $C$}{R, C})}

The computation proceeds in three phases.

\paragraph{Phase 1 --- Outer join preservation.}
Entries from preserved sides are copied unconditionally:
\begin{align*}
  R'_{\mathrm{outer}} &=
    \bigl(\Rf \text{ if } \outerf \text{ else } \emptyset\bigr)
    \;\cup\;
    \bigl(\Rp \text{ if } \outerp \text{ else } \emptyset\bigr)
  \\
  C'_{\mathrm{outer}} &=
    \bigl(\Cf \text{ if } \outerf \text{ else } \emptyset\bigr)
    \;\cup\;
    \bigl(\Cp \text{ if } \outerp \text{ else } \emptyset\bigr)
\end{align*}

\paragraph{Phase 2 --- Guard.}
If any of the following conditions fail, set
$R' = R'_{\mathrm{outer}}$, $C' = C'_{\mathrm{outer}}$, and stop:
\[
  \nn \;\land\; f \in \Nf \;\land\; p \in \Rp
\]

The guard requires three things:
\begin{itemize}[nosep]
\item $\nn$: the FK columns carry \textsc{not null} constraints on the
  base table.
\item $f \in \Nf$: the FK columns have not been corrupted by
  null-padding from prior outer joins.  Even with \textsc{not null}
  base constraints, prior joins may have introduced ghost rows with
  \textsc{null} FK values; $f \in \Nf$ confirms this has not happened.
\item $p \in \Rp$: the referenced base table's complete row set
  appears in $\Jp$'s result.
\end{itemize}
Without all three, a referencing row may fail to find its match:
either because its FK value is \textsc{null} or corrupted, or because
the matching referenced row is missing.

\paragraph{Phase 3 --- Chain extension.}
When the guard passes:

\medskip\noindent
\emph{Step 3a.}\; Compute the anchor set---tables that reach~$f$:
\[
  \anchor = \set{A : (A, f) \in \Cf}
    \;\cup\; \bigl(\set{f} \text{ if } f \in \Rf\bigr)
\]
Each member of $\anchor$ is row-preserving and connected to~$f$
through a NOT NULL FK chain (or is $f$ itself when $f \in \Rf$).

\medskip\noindent
\emph{Step 3b.}\; Inherit from the referencing side.  Copy $R$ and $C$
entries rooted at anchor members:
\begin{align*}
  R'_{\mathrm{inherit}} &= \anchor \\
  C'_{\mathrm{inherit}} &=
    \set{(A, Y) \in \Cf : A \in \anchor}
\end{align*}
Every member of $\anchor$ is in $\Rf$ (by the chain invariant and
the condition on $f$ in Step~3a), so
$R'_{\mathrm{inherit}} \subseteq \Rf$.  Tables in $\Rf$ that do not
reach~$f$ are excluded because the join may drop some of their rows.

\medskip\noindent
\emph{Step 3c.}\; Extend chains across the FK boundary.  Each anchor
member now reaches~$p$, and transitively everything $p$ reaches:
\[
  C'_{\mathrm{extend}} =
    \set{(A, p) : A \in \anchor}
    \;\cup\;
    \set{(A, Y) : A \in \anchor,\; (p, Y) \in \Cp}
\]

\paragraph{Final result:}
\begin{align*}
  R' &= R'_{\mathrm{outer}} \cup R'_{\mathrm{inherit}} \\
  C' &= C'_{\mathrm{outer}} \cup C'_{\mathrm{inherit}}
       \cup C'_{\mathrm{extend}}
\end{align*}


\subsection{GROUP BY}

When a derived relation applies GROUP BY, it modifies the four
properties of the underlying relation's output before the result is
used in further joins.  Let $(U_{\mathrm{in}}, R_{\mathrm{in}},
N_{\mathrm{in}}, C_{\mathrm{in}})$ be the properties of the
underlying relation.

\paragraph{Conditions for uniqueness restoration.}
GROUP BY can restore uniqueness when:
\begin{itemize}[nosep]
\item Every grouping expression is a direct column reference (no
  computed expressions)
\item All grouping columns reference the same base table~$t$
\item The grouping columns cover all columns of some unique key on~$t$
\end{itemize}

\paragraph{Effect on the four properties.}
When the conditions above are met:
\begin{itemize}
\item $U$: add $t$ to $U_{\mathrm{in}}$ (uniqueness restored).  This
  holds even when a group-level filter is present, since uniqueness is
  about deduplication, not row completeness.
\item $R$: propagate $R_{\mathrm{in}}$, but only if no group-level
  filter is present (a group-level filter may discard groups, breaking
  row preservation).
\item $C$: propagate $C_{\mathrm{in}}$ (same condition as $R$).
\item $N$: propagate $N_{\mathrm{in}}$ only if all columns of the
  matched unique key carry \textsc{not null} constraints.  If any
  column is nullable, GROUP BY collapses all \textsc{null} values in
  that column into a single group, producing a \textsc{null} key value
  that may not correspond to any base table row.  In this case, $N$ is
  set to~$\emptyset$.
\end{itemize}

When GROUP BY does not restore uniqueness, or when a group-level
filter is present, $R$, $C$, and $N$ are not propagated.


%% ====================================================================
\section{Acceptance Criterion}

At the entry point, the algorithm computes the four properties for the
referenced subtree~$\Jp$.  The FK join is accepted if and only if:
\[
  p \in U \;\land\; p \in R \;\land\; p \in N
\]
where $U$, $R$, $N$ are the output properties of~$\Jp$.

\begin{itemize}
\item $p \in U$: the referenced base table's unique key is preserved,
  ensuring each referencing row matches at most one referenced row.
\item $p \in R$: the referenced base table's complete row set appears,
  ensuring every valid FK value finds its match.
\item $p \in N$: the referenced base table's key-column values are
  intact, ensuring the unique key functions correctly (no spurious
  \textsc{null}s from outer-join padding or GROUP BY on nullable
  columns).
\end{itemize}

For base tables accessed directly (not through a derived relation),
all three conditions hold trivially.


%% ====================================================================
\section{Soundness}

\begin{theorem}
If the acceptance criterion holds ($p \in U \land p \in R \land p \in N$
on the referenced side's output), then
$\pi_{\mathbf{b}}(S) = \pi_{\mathbf{b'}}(s)$.
\end{theorem}

\begin{proof}[Proof sketch]
The three conditions together establish that the referenced side~$S$
contains the same key values as the base table~$s$ with the same
multiplicities:
\begin{itemize}
\item $p \in R$ guarantees that all rows of $s$ appear in $S$ (no row
  loss).  Every key value in $\pi_{\mathbf{b'}}(s)$ appears in
  $\pi_{\mathbf{b}}(S)$ with at least its original multiplicity.
\item $p \in U$ guarantees that $p$'s unique key is preserved through
  all joins in $\Jp$ (no row duplication).  Since the key columns
  of~$s$ carry a unique constraint, and this uniqueness is preserved,
  no key value in $\pi_{\mathbf{b}}(S)$ appears with multiplicity
  greater than one.
\item $p \in N$ guarantees that the key-column values have not been
  corrupted by \textsc{null}s.  Without this, ghost rows from outer
  joins could introduce \textsc{null} key values, or GROUP BY could
  collapse multiple \textsc{null} rows into one, violating the
  multiset equality.
\end{itemize}
Together, these imply $\pi_{\mathbf{b}}(S) = \pi_{\mathbf{b'}}(s)$:
every key value appears exactly once, and no spurious values are
introduced.
\end{proof}


\subsection{Chain Invariant Proof}

The chain invariant (\S5) is proved by structural induction on the
join tree.

\paragraph{Base case.}
At a leaf node, $C = \emptyset$ and the invariant holds vacuously.

\paragraph{Inductive step.}
Assume the invariant holds for both subtrees $\Jf$ and $\Jp$.  We
show it holds for $J = (\Jf \fkjoin{f \to p,\;\nn,\;\tau} \Jp)$.

\medskip\noindent
\emph{Outer-preserved entries.}\;
When a side is preserved by the outer join, all its rows appear in the
result (possibly padded with \textsc{null}s on the other side).  The
$R$ and $C$ entries from that side remain valid because no rows are
lost.  The invariant is maintained by direct copy.

\medskip\noindent
\emph{Guard condition.}\;
When $\lnot\nn \lor f \notin \Nf \lor p \notin \Rp$, no new chain is
valid: either \textsc{null} FK values could cause row loss (violating
condition~1 of the invariant), or the FK columns have been corrupted
by prior outer joins (also violating condition~1), or the referenced
table is incomplete (violating condition~3).  Returning with only
outer-preserved sets is correct.

\begin{lemma}[Key lemma]
If $A \in \anchor$, then every row of $A$ in $\Jf$'s result survives
the join.
\end{lemma}

\begin{proof}
$A \in \anchor$ means either $A = f$ with $f \in \Rf$, or
$(A, f) \in \Cf$.  In both cases, $A$ has a NOT NULL FK chain ending
at~$f$, so every row of~$A$ in $\Jf$ is associated with a specific
value of $f$'s FK columns.  Since $f \in \Nf$ (from the guard), these
FK values have not been corrupted by prior joins.  Since $\nn$ holds,
they are non-null.  Since $p \in \Rp$, the referenced table is
complete.  The FK constraint guarantees that every non-null FK value
references an existing row.  Therefore every row of~$A$ matches
exactly one row in $\Jp$ and survives the join.
\end{proof}

\noindent
\emph{Inherit step.}\;
Chains $(A, Y)$ from $\Cf$ are copied only when $A \in \anchor$.  By
the lemma, $A$ remains row-preserving.  By the inductive hypothesis,
the chain from $A$ to $Y$ was valid in~$\Jf$.  Since $A$'s rows are
preserved, the chain remains valid in~$J$.

\medskip\noindent
\emph{Extend step.}\;
The new chain $(A, p)$ is valid: $A \in \anchor$ ensures $A \in R'$,
the edge $f \to p$ is a NOT NULL FK join, and $p \in \Rp$ holds at
the time of processing---satisfying all three conditions of the chain
invariant.  Extensions $(A, Y)$ for $(p, Y) \in \Cp$ compose the path
from $A$ through $p$ to $Y$; by the inductive hypothesis on~$\Jp$,
the chain $p \to \cdots \to Y$ was valid, and prepending the path from
$A$ through $f$ to $p$ preserves validity.


\subsection{Null Preservation Correctness}

\emph{Base case.}\;
$N = \set{t}$ is trivially correct: a single base table has no joins,
so its key-column values are intact.

\medskip\noindent
\emph{Remove step (outer joins).}\;
When a side is outer-preserved, unmatched rows from the preserved side
produce ghost rows with all-\textsc{null} columns on the inner side.
Removing the inner side's $N$ entries is correct because those tables'
key columns may now contain spurious \textsc{null}s.

Exception: when the referencing side is outer-preserved and $\nn$
holds, every referencing row has a non-null FK value and the FK
constraint guarantees a match.  No unmatched referencing rows exist,
so no ghost rows appear on the referenced side.  Therefore $\Np$ is
not removed.

\medskip\noindent
\emph{Remove step (GROUP BY).}\;
When GROUP BY on a nullable column restores uniqueness, it collapses
all \textsc{null} values into a single group.  This group's key value
is \textsc{null}, which may not correspond to any base table row (the
base table may have multiple rows with \textsc{null} in that column,
or none).  Removing the table from $N$ is correct.


\subsection{Selective Inheritance}

Tables in $\Rf$ that are not in $\anchor$ are correctly excluded from
$R'$.  Such a table~$B$ has no NOT NULL FK chain to~$f$.  The join on
$f$'s FK columns may drop some of $B$'s rows---specifically, rows
paired with $f$-rows whose FK values are \textsc{null} or whose FK
columns have been corrupted by prior joins.  Without a guarantee that
all of $B$'s rows survive, $B$ cannot be in~$R'$.


%% ====================================================================
\section{Conservativeness}

The analysis is deliberately conservative.  The following classes of
valid queries are rejected:

\begin{enumerate}
\item \textbf{Non-FK joins preserving rows.}\;
  A non-FK join inside a derived relation may happen to preserve all
  rows (e.g., a join whose condition matches every row), but the
  analysis cannot determine this.  Any non-FK join causes all tracked
  properties to be lost.

\item \textbf{Set operations.}\;
  Operations that combine multiple relations (union, intersection,
  difference) are not analyzed, even when they might preserve the
  required properties.

\item \textbf{Complex grouping.}\;
  The analysis requires grouping columns to be direct column
  references from a single base table covering a unique key.  It
  rejects grouping on expressions (even deterministic ones), grouping
  on columns from multiple tables, and grouping on columns from nested
  derived relations.

\item \textbf{Deduplication.}\;
  Even when deduplication would not change the result (e.g., applied
  to an already-unique relation), it is rejected because the analysis
  cannot verify this statically.

\item \textbf{Semantically vacuous filters.}\;
  A filter that happens to exclude no rows (e.g., a tautological
  predicate) breaks row preservation in the analysis, even though the
  data is unchanged.

\item \textbf{Transitive reasoning across non-FK joins.}\;
  If two FK chains are connected by a non-FK join that happens to
  preserve rows, the analysis cannot chain through the non-FK join.
\end{enumerate}

This conservativeness is by design: false negatives (rejecting valid
joins) are acceptable; false positives (accepting invalid joins) are
not.

\end{document}
